% !TEX root = ../CerdaMS.tex

\section{Preliminaries}

We recall some basic notions and results concerning the \lam-calculus, and its fragment without weakening known as the \lamI-calculus. For more information, we refer to Barendregt's book \cite{Bare}.

\subsection{The \lam-calculus, in a nutshell} 
Let us fix an infinite set $\Var$ of variables. The set $\Lam$ of 
\emph{\lam-terms} is defined by induction
\bnf[\Lam]{M,N}{ x \mid \lam x.M \mid MN }[\text{for $x \in \Var$}]
We use parentheses when needed.
We assume that application associates on the left, 
and has higher precedence than abstraction.
E.g., $\lam x.\lam y.\lam z.xyz$ stands for $\lam x.(\lam y.(\lam z.(xy)z))$.
We write $\lam x_1 \ldots x_n.M$ or even $\lam\vec x.M$ 
for $\lam x_1.\ldots \lam x_n.M$, 
and $N^n(M)$ for $N(N(\cdots N(M)\cdots))$, where $N$ occurs $n$ times. 
In particular, when $n = 0$, we get $\lam\vec x.M = N^n(M) = M$.

The set $\fv(M)$ of \emph{free variables of $M$} is defined as usual, and we say that $M$ is \emph{closed} whenever $\fv(M) = \emptyset$.
Hereafter, \lam-terms are considered modulo \emph{$\alpha$-conversion}
(see \cite[\S 2.1]{Bare}).

\begin{example}\label{ex:combinators}
The \lam-terms below are used throughout the paper to construct examples:
\[
	\comb{I} \eqdef \lam x.x,\quad 
	 \comb{K} \eqdef \lam xy.x,\quad 
	 \comb{F} \eqdef \lam xy.y,\quad 
	 \comb{B} \eqdef \lam fgx.f(g(x)),\quad
 	 \comb{D} \eqdef \lam x.xx,\quad
	 \Om \eqdef \comb{DD}.
\]
The \lam-term $\comb{I}$ is the identity, $\comb{K,F}$ are the projections, $\comb{B}$ is the composition, $\comb{D}$ is the diagonal and $\Om$ a looping term.
%Given $n\in\Nat$, the $n$-th \emph{Church numeral} is defined by $\comb{c}_n = \lam sz.s^n(z)$.
The \emph{pairing of $M,N$} is encoded as $[M,N] \eqdef \lam x.xMN$, for $x$ fresh. 
\end{example}
The \emph{$\beta$-reduction} $\to_\beta$ is generated by the rule
$
	(\lam x.M)N\to M\subst{N}
$, where $(\lam x.M)N$ is called a \emph{$\beta$-redex} and $M\subst{N}$ stands for the \emph{capture-free substitution} of $N$ for all free occurrences of $x$ in $M$. 
We let $\redd[\beta]$ (resp. $=_\beta$) be the reflexive-transitive (and symmetric) closure of $\to_\beta$.
We say that $M$ \emph{is in $\beta$-normal form ($\beta$-nf)} if $M$ does not contain any $\beta$-redex, and that $M$ \emph{has a $\beta$-nf} if $M\redd[\beta] N$, for some $N$ in $\beta$-nf.
Since $\red[\beta]$ is confluent, such $N$ must be unique.

\begin{remark}\label{rem:hnfs}
Any \lam-term $M$ can be written in one of the following forms: either $M = \lam\vec{x}.yM_1\cdots M_k$, in which case $M$ is called a \emph{head normal form} (hnf) and $y$ its \emph{head variable}; or $M = \lam\vec{x}.(\lam y.P)QM_1\cdots M_k$, where $(\lam y.P)Q$ is referred to as its \emph{head redex}.
\end{remark}
%\comasi{Ici, je pense qu'il pourrait être utile d'ajouter que les formes normales sont les termes héréditairement en hnf.}
The \emph{head reduction} $\red[h]$ is the restriction of $\red[\beta]$ obtained by contracting the head redex.

\begin{definition} A \lam-term $Y$ is a \emph{fixed-point combinator} (fpc) if it satisfies $Y x =_\beta x(Yx)$.
\end{definition}
Notice that we do not require that fixed-point combinators are closed \lam-terms.
We recall Curry's combinator $\comb{Y}$ and define, for all $l\in\Var$, 
Polonsky's $\Theta_l$ and Klop's $\Bible$ \cite{ManzonettoPSS19}:
\begin{equation}\label{eq:fpc}
	\comb{D}_f \eqdef \lam x.f(xx),\quad
	\comb{Y} \eqdef \lam f.\comb{D}_f\comb{D}_f,\quad
	\comb{V} \eqdef \lam vlf.f(vvlf),\quad
	\Theta_l \eqdef \comb{VV}l,\quad
	\Bible \eqdef \lam e.\comb{BYB}el.
\end{equation}
It is easy to check that $\comb{Y}$ and $\Theta_l$ are fpcs.
Regarding the \emph{Bible} $\Bible$, which owes its name to its distinctive spelling, we get $\Bible x =_\beta \comb{BYB}xl =_\beta \comb{Y}(\comb{B}x)l =_\beta \comb{B}x(\comb{Y}(\comb{B}x))l =_\beta x(\comb{Y}(\comb{B}x)l) =_\beta x(\Bible x)$.
Note that the variable $l$ remains in a passive position along the reduction of $\Bible$ (and of $\Theta_l$), therefore $\Bible[N]:= \Bible\subst[l]{N}$ (resp.\ $\Theta_N:= \Theta_l\subst[l]{N}$) remains a fpc, for all $N\in\Lambda$.

%By confluence, $Y$ is an fpc if $Y\redd[h] \lam f.fY'$, for some $Y'$ such that $\lam f.Y'$ is itself an fpc. So we can continue the head reduction within $Y'\redd[h] fY''$, and so forth.
Although fpcs $Y$ are not $\beta$-normalizable, they differ from $\Om$ since they produce increasing stable amounts of information along reduction: $Y\redd[\beta] \lam f.fY_1\redd[\beta] \lam f.f^n(Y_n)\redd[\beta]\cdots$.
Barendregt introduced a notion of `evaluation tree' to capture this infinitary behaviour~\cite{Barendregt77}.\par\addvspace{1.5ex}

\noindent\begin{minipage}{\textwidth}
\begin{definition}\label{def:BT} 
The \emph{Böhm tree} $\BT{M}$ of a \lam-term $M$ is coinductively defined as follows.
\begin{enumerate}
\item If $M$ has a hnf, that is $M \redd[h] \lam x_1\dots x_n.yM_1\cdots M_k$, then\\
\[
\begin{tikzpicture}
\node 		(BTM) at (-23pt,0)					{$\BT M \eqdef$};
\node[n] 	(y) 	at (1.1,0) 		{$\lam x_1\dots x_n.\,y$};
\node[n] 	(BTM1) 	at (.8,-.8) 	{$\BT{M_1}$};
\node 				at (2,-.8) 		{\dots};
\node[n] 	(BTM2) 	at (3.2,-.8) 	{$\BT{M_k}$};
\draw[b] ($(y.south east)+(-8pt,0)$) -- (BTM1.north);
\draw[b] (y.south east) -- (BTM2.north);
\end{tikzpicture}
\]
\item $\BT{M} \eqdef \bot$, otherwise. The constant $\bot$ represents the absolute lack of information.
\end{enumerate}
\end{definition}
\end{minipage}

\begin{example}\label{ex:BTs} 
Using the terms from Example~\ref{ex:combinators} and \eqref{eq:fpc}, let us construct some Böhm trees.
\bsub
\item If $M$ has a $\beta$-normal form $N$ then, up to isomorphism, $\BT{M} = N$. E.g., $\BT{\comb{DI}} = \comb{I}$.
\item Since $\Om$ and $\comb{YK}$ have no hnf, we get $\BT{\Om} = \BT{\comb{YK}} = \bot$.
\item If $Y$ is an fpc, then $\BT{Y} = \lam f.f(f(f(\cdots)))$. Thus, $\BT{\comb{Y}} = \BT{\Bible} = \BT{\Theta_l}$.
\item\label{ex:BTs4} 
	Using the fpc $\comb{Y}$, one can define \lam-terms $\comb{E}_x,\comb{O}_x$ satisfying
	$\comb{E}_x =_\beta {[}\Om, [\Om x, \comb{E}_x]]$ and 
	$\comb{O}_x =_\beta {[}\Om x, [\Om, \comb{O}_x]]$. 
These two \lam-terms both construct a `stream' $[M_1,[M_2,[M_3,[\cdots]]]]$, but the former places $\Om x$ at even positions and $\Om$ at odd ones, while the latter does the converse.
This difference is however forgotten in their Böhm trees:
\[
	\BT{\comb{E}_x} = \BT{\comb{O}_x} = [\bot,[\bot,[\bot,[\bot,[\,\cdots]]]]]\hfill\textrm{ (inline depiction of a tree)}
\]
\item Check that, given $\comb{R} = \lam ylx.x(y(xl))$, we have $\BT{\comb{YR}l} = \lam x_0.x_0(\lam x_1.x_1(\lam x_2.x_2(\cdots)))$.
\esub
\end{example}

The B\"ohm tree equality induces an equivalence on \lam-terms: $M =_\cB N$ iff $\BT{M} = \BT{N}$. 
As shown in \cite{Levy76,Wadsworth78}, $\cB$ is a \emph{\lam-theory} because it contains $=_\beta$ and is \emph{compatible} with abstraction and application: $M =_\cB M'$ entails $\lam x.M =_\cB \lam x.M'$, $MN =_\cB M'N$ and $NM =_\cB NM'$.

\subsection{The \lamI-calculus}\label{ssec:lamI}

The set $\LamI$ of \lamI-terms is the subset of $\Lam$ consisting of those \lam-terms 
in which every abstraction binds \emph{at least} one occurrence of the abstracted variable.
Note that  $\red[\beta]$ induces a reduction on $\LamI$ since $M\in\LamI$ and $M \red[\beta] N$ entail $N\in\LamI$. 
In other words, the property of being a \lamI-term is preserved by $\red[\beta]$.
Similarly, the \lamI-calculus inherits from \lam-calculus all the notions previously defined like head reduction, Böhm trees, etc.
For convenience, we consider an alternative definition of $\LamI$ simultaneously defining $\lamI$-terms and their sets of free variables.

\begin{definition}\label{def:LamI}\mbox{}
\bsub
\item For all $X \subseteq \Var$, $\LamI(X)$ is defined as 
	the smallest subset of $\Lam$ such that:
\[
	\infer{x\in\LamI(\set{x})}{}\qquad
	\infer{\lam x.M\in\LamI(X\setminus\set{x})}{M\in\LamI(X)\quad x\in X}\qquad
	%\infer{\lam x.M\in\LamI(X)}{M\in\LamI(\set{x}\cup X)\quad x\notin X}\qquad
	\infer{MN\in\LamI(X\cup Y)}{M\in\LamI(X) & N\in\LamI(Y)}
\]
%The above definition differs from the usual definition of (regular) 
%\lam-terms with their set of free variables only in the side condition of 
%the second clause (namely, `$x\in X$').
\item Finally, the set of all \lamI-terms is given by 
	the disjoint union
	$\LamI = \bigdunion_{X\subseteq\Var}\LamI(X)$.
%$\LamI = \set{ M\in\Lam \mid \exists X\subseteq\Var\,.\, M\in\LamI(X)}$.
\esub
\end{definition}

\begin{remark}\label{rem:aboutLambdaI}\mbox{}
\bsub
\item Note that $M \in \LamI(X)$ means that $M$ is a \lamI-term such that $\fv{(M)}=X$.
As a consequence, if $M\in\LamI(X)$ then the set $X$ must be finite.
\item\label{rem:aboutLambdaI2} 
	Each set $\LamI(X)$ is closed under $\beta$-conversion.
\item Our examples $\comb{I},\comb{B},\comb{D},\Om,\comb{Y},\Bible,\Theta_l,\comb{E}_x,\comb{O}_x,\comb{YR}l$ are \lamI-terms, while $\comb{K},\comb{F},\comb{YK}\notin\LamI$.
\esub
\end{remark}

\begin{definition}
A \emph{\lamI-theory} $\cT$ is given by an equivalence $\mathord{=_\cT} \subseteq \LamI^2$ containing $\beta$-conversion and compatible with application and abstraction. 
In the context of \lamI-calculus, 
the compatibility with abstraction becomes: 
$M =_\cT M'$ and $x\in\fv(M)\cap\fv(M')$ 
imply $\lam x.M =_\cT \lam x.M'$.
\end{definition}
The only subtlety in the definition above stands in the side condition on the compatibility with abstraction.
However, this small difference has significant consequences: e.g., the theory axiomatized by equating all terms without a $\beta$-nf is consistent as a \lamI-theory, but inconsistent as a \lam-theory.
The following theorem collects the main properties of the \lamI-calculus.
\begin{theorem}\label{thm:LamIMainProps}\mbox{}
\bsub
\item\label{thm:LamIMainProps1} 
	Every computable numerical function is \lamI-definable~\cite[Thm.~9.2.16]{Bare}.
\item\label{thm:LamIMainProps2} 
	Strong and weak $\beta$-normalization coincide for the \lamI-calculus~\cite[Thm.~11.3.4]{Bare}.
\item\label{thm:LamIMainProps3} 
	Every \lam-theory is also a \lamI-theory, while the converse does not hold. 
\esub
\end{theorem}
%Some comments on these results.
%To prove \eqref{thm:LamIMainProps1} (aka.~Turing completeness) it is sufficient to redefine Church's numeral zero $\comb{c}_0^\comb{I} = \lam xy.x\comb{II}y\in\LamI$ (see ). 
%Consequently, the \lam-terms $\comb{E}_x,\comb{O}_x$ from Example~\ref{ex:BTs}\eqref{ex:BTs4} can be defined in such a way that they belong to $\LamI$.


By \eqref{thm:LamIMainProps3}, Böhm trees induce a \lamI-theory, but we argue that this theory does not respect the spirit of $\LamI$. 
E.g., the variable $l$ is never erased along the reduction of, say, $\Bible$, but it disappears in its B\"ohm tree. Using the terminology of \cite{Bare}, one says that $l$ is \emph{pushed to infinity} in $\BT{\Bible}$.
This shows that the Böhm tree of the \lamI-term $\lam l.\Bible\in\LamI$ is not a $\lamI$-tree: the outer \lam-abstraction `$\lam l$' does not bind any free occurrence of $l$ in $\BT{\Bible} = \lam f.f(f(f(\,\cdots)))$.
